\chapter{Виртуальный цветной мост и граница дополнения Шура}
\label{ch:v7-virtual-colored-bridge-schur-complement-gate}

\section{Какая переменная действительно интегрируется}

Предыдущий гейт оставил возможность виртуального участия двух цветных
рёбер при нулевых одноточечных средних. Однако формула дополнения Шура
имеет два разных физических чтения. В гейтах messenger-цепей Тома IV
исключался тяжёлый подблок линейного фермионного оператора. Здесь цветные
рёбра сами являются бозонными полями. Поэтому классическое исключение и
детерминант их флуктуаций должны проверяться раздельно.

Обозначим
\begin{equation}
 a=z_{u_RX_L}\in\overline{\mathbf3},\qquad
 b=z_{Y_RQ_L}\in\mathbf3,
 \qquad
 p=z_{X_Le_R},\qquad q=z_{L_LY_R}.
 \label{eq:v7-virtual-bridge-fields}
\end{equation}
После фиксации двух старых корней укоренённый цикл пропорционален
$pq\,a\cdot b+\mathrm{c.c.}$ Минимальная тяжёлая квадратичная форма на
каждой цветовой компоненте задаётся эрмитовым блоком
\begin{equation}
 K(pq)=
 \begin{pmatrix}
  M_a^2&-\kappa\overline{pq}\\
  -\kappa pq&M_b^2
 \end{pmatrix},
 \qquad
 \Delta(pq)=\det K=M_a^2M_b^2-\kappa^2|pq|^2.
 \label{eq:v7-virtual-bridge-heavy-block}
\end{equation}
Цветосохраняющая тяжёлая область требует
\begin{equation}
 \Delta(pq)>0.
 \label{eq:v7-virtual-bridge-positive-domain}
\end{equation}

\section{Классическое исключение ничего не индуцирует}

В текущем цикле отсутствует член, линейный по $a$ или $b$. Поэтому внутри
области \eqref{eq:v7-virtual-bridge-positive-domain} тяжёлые стационарные
уравнения имеют единственное решение
\begin{equation}
 a=b=0.
 \label{eq:v7-virtual-bridge-tree-solution}
\end{equation}
Подстановка этого решения возвращает исходный лёгкий потенциал без
циклической поправки. Обычная формула
$D_{\rm eff}=D_{ll}-D_{lh}D_{hh}^{-1}D_{hl}$ здесь неприменима без
дополнительного линейного источника: она относится к исключению состояний
линейного оператора, а не к автоматическому исключению его полевых
коэффициентов.

\section{Конечномерный детерминант}

Для трёх комплексных цветовых копий гауссов интеграл даёт
\begin{equation}
 \Gamma_{0}(p,q)
 =3\log\Delta(pq)+\mathrm{const}.
 \label{eq:v7-virtual-bridge-zero-mode-logdet}
\end{equation}
Разложение при малом $|pq|$ равно
\begin{equation}
 \Gamma_{0}(p,q)
 =3\log(M_a^2M_b^2)
 -\frac{3\kappa^2}{M_a^2M_b^2}|pq|^2
 +O(|pq|^4).
 \label{eq:v7-virtual-bridge-quartic-expansion}
\end{equation}
Первый индуцированный оператор является gauge-инвариантным и имеет
четвёртую степень по лёгким полям. Следовательно,
\begin{equation}
 \Hess_{p=q=0}\Gamma_0=0.
 \label{eq:v7-virtual-bridge-zero-light-hessian}
\end{equation}
Виртуальная цветная пара создаёт нелинейное притяжение между $p$ и $q$,
но не запускает их из полного нуля.

Собственные значения $K$ равны
\begin{equation}
 \lambda_{\pm}
 =\frac12\left[M_a^2+M_b^2
 \pm\sqrt{(M_a^2-M_b^2)^2+4\kappa^2|pq|^2}\right]
 \label{eq:v7-virtual-bridge-heavy-eigenvalues}
\end{equation}
с цветовой кратностью три. При $\Delta\to0$ значение $\lambda_-$ стремится
к нулю: интегрируемая мода перестаёт быть тяжёлой именно на границе
сильного конечномерного эффекта. Это повторяет предупреждение детерминантного
гейта моста Тома VI, хотя полный четырёхмерный предел требует отдельного
анализа.

\section{Четырёхмерный след логарифма}

После добавления канонической кинетики первый член разложения имеет вид
\begin{equation}
 \Gamma_{4}^{(1)}
 =-3\kappa^2|pq|^2
 \int\frac{d^4k}{(2\pi)^4}
 \frac{1}{(k^2+M_a^2)(k^2+M_b^2)}.
 \label{eq:v7-virtual-bridge-four-dimensional-term}
\end{equation}
Радиальный ультрафиолетовый хвост пропорционален $dk/k$. Поэтому это
логарифмически перенормируемая квартичная вершина, а не конечный
коэффициент, однозначно заданный конечным графом. Необходимо фиксировать
контрчлен, масштаб согласования и полный gauge/BV-набор. Такое разделение
tree-level и one-loop вкладов стандартно при интегрировании тяжёлых полей
методом фонового поля
\cite{v7-virtual-bridge-hlm,v7-virtual-bridge-dittmaier}.

\begin{theorem}[Нелинейный проход без самостоятельного запуска]
В положительной тяжёлой области виртуальная цветная пара имеет нулевые
классические средние и сохраняет $SU(3)_c$. Её классическое исключение не
создаёт лёгкого оператора. Конечномерный детерминант индуцирует отрицательный
gauge-инвариантный член $|pq|^2$, но его лёгкий гессиан в полном нуле равен
нулю. В четырёх измерениях коэффициент этого члена требует
перенормировочного согласования. Следовательно, виртуальный мост допустим
как нелинейная связь, но не как самостоятельный coefficient-free родитель
цветосохраняющего запуска.
\end{theorem}

\begin{nogocriterion}[Запрет подмены двух дополнений Шура]
Нельзя переносить messenger-формулу исключения тяжёлых фермионных состояний
на бозонные коэффициенты рёбер без явного линейного источника. Нельзя также
считать коэффициент $-3\kappa^2/(M_a^2M_b^2)$ физически предсказанным до
вывода $M_a,M_b,\kappa$, кинетической меры и четырёхмерного правила
перенормировки из одного действия.
\end{nogocriterion}

\begin{protocol}
\begin{enumerate}
 \item тяжёлых цветных полей: \textbf{два, $\overline3+3$};
 \item цветовая кратность блока: \textbf{$3$};
 \item область сохранения тяжёлой щели: \textbf{$\Delta>0$};
 \item классическое решение в этой области: \textbf{$a=b=0$};
 \item tree-level поправка лёгкого потенциала: \textbf{$0$};
 \item первый конечномерный determinant-оператор:
 \textbf{$-3\kappa^2|pq|^2/(M_a^2M_b^2)$};
 \item его степень по лёгким полям: \textbf{$4$};
 \item лёгкий гессиан в полном нуле: \textbf{$0$};
 \item самостоятельный цветосохраняющий запуск: \textbf{нет};
 \item нелинейная singlet-связь: \textbf{да};
 \item закрытие тяжёлой щели при $\Delta=0$: \textbf{да};
 \item четырёхмерный коэффициент конечен без контрчлена: \textbf{нет};
 \item единая нормировка $M_a,M_b,\kappa$: \textbf{не выведена};
 \item итог: \textbf{частичный структурный проход, родительский no-go};
 \item следующий гейт: \textbf{происхождение цветосохраняющего
 квадратичного селектора}.
\end{enumerate}
\end{protocol}

Машинный сертификат сохранён в
\path{s2t_v7_virtual_colored_bridge_schur_complement_gate_results.json}.

\begin{thebibliography}{9}
\bibitem{v7-virtual-bridge-hlm}
B.~Henning, X.~Lu, H.~Murayama,
\emph{One-loop Matching and Running with Covariant Derivative Expansion},
Journal of High Energy Physics 01 (2018), 123; arXiv:1604.01019.
\bibitem{v7-virtual-bridge-dittmaier}
S.~Dittmaier, S.~Schuhmacher, M.~Stahlhofen,
\emph{Integrating out heavy fields in the path integral using the
background-field method}, European Physical Journal C 81 (2021), 826;
arXiv:2102.12020.
\end{thebibliography}